\section{Related work}\label{sec:related}

\subsection{Efficient reasoning models}
LRMs allocate additional test-time computation to explicit reasoning traces often with specific delimiters such as \texttt{<think>}\,\ldots\,\texttt{</think>}~\cite{guo2025deepseek, jaech2024openai}. This \emph{think-before-answer} format enables iterative exploration and self-verification, but it is frequently applied \emph{uniformly} across inputs. Such uniform long-form reasoning leads to the \emph{overthinking} challenge with high latency and compute on intrinsically easy queries and sometimes producing redundant reasoning tokens \cite{feng2025efficient, sui2025stop}.

A broad line of work has pursued efficiency via \emph{adaptive computation} and \emph{early stopping}. Classical dynamic-computation methods and early-exit approaches terminate generation once confidence criteria are met, reducing unnecessary computation without changing the core model objective \cite{schuster2022confident, bae2025mixture}. Recent reasoning-focused extensions similarly aim to avoid excessive reasoning tokens by learning or heuristically determining when further context is unlikely to help, and by truncating or suppressing low-information reasoning segments \cite{sui2025stop}.

Another direction controls reasoning \emph{length} through prompt- or budget-based strategies. Methods that rely on fixed or predicted token budgets can reduce verbose CoT while maintaining accuracy, but often require external budget selection rules or calibrated estimators, which may not transfer cleanly across domains \cite{han2024token}. In parallel, internal-rationale approaches such as Quiet-STaR \cite{zelikman2024quiet} encourage models to form latent rationales during generation, while simplified recipes such as s1 \cite{muennighoff2025s1} show that carefully constructed 1k data and prompting can recover much of test-time scaling benefits with minimal training overhead.

As an orthogonal axis over these approaches, our goal is not to enforce a budget or modify the architecture, but to \emph{internalize the gating decision itself} whether to open a long explicit reasoning trace in the first place. \Ours~optimizes only the \emph{root token} that triggers the \texttt{<think>} branch, enabling adaptive reasoning with minimal training footprint and without backpropagating over whole reasoning trajectories.

\subsection{Difficulty-aware policy optimization}
The question of \textit{learning when to reason} has recently been reframed as \emph{difficulty-aware policy optimization}: the model should allocate explicit reasoning only when needed, while responding directly to easy or non-reasoning queries. This direction is especially relevant under the explicit \texttt{<think>} styled format, where the core idea is a discrete \emph{branching action} that enters a long reasoning mode versus produces a direct response \cite{Zhang2025AdaptThink, Bao2025Thinkless, wu2025arm, wang2025adareasoner}.

Representative methods learn such policies through RL or preference learning. AdaptThink \cite{Zhang2025AdaptThink} trains a policy to select \emph{thinking} versus \emph{non-thinking} modes conditioned on difficulty, while Thinkless \cite{Bao2025Thinkless} uses preference/RL signals to reduce redundant thinking tokens in general conversations. Beyond binary gating, ARM \cite{wu2025arm} expands the action space to multiple reasoning formats (e.g., direct answer, short CoT, long CoT, code), and AdaReasoner \cite{wang2025adareasoner} treats reasoning configurations as actions in a model-agnostic RL framework. Together, these works show that adaptive reasoning is learnable and can outperform static long-CoT inference.

However, existing policy optimization approaches face practical hurdles that motivate our design. 
First, many studies are benchmark-centric and largely evaluated on math-heavy suites, leaving open how well their learned routing \emph{calibrates} (e.g., whether the thinking rate remains aligned with difficulty on open-domain and scientific queries). 
Second, optimizing policies over whole long reasoning trajectories is often \emph{training-cost dominated}: GRPO-style post-training incurs expensive cost of VRAM, gradient computation, and  wall-clock costs \cite{guo2025deepseek}, also detailed in \S\ref{subsec:rpo} and \Autoref{tab:pflops_adaptive_reasoning_8b}. 
Third, effective training frequently depends on curated external prompts, data from larger teacher models, paired supervision, or proprietary query streams for routing signals, and these factors are rarely opened to the community. 

